\subsection{Complete display of the seventeen-variable equation}

For an explicit standalone specification, the seventeen unrestricted
integer coordinates are
\[
 J_3,C,z,f,h,r,v,w_3,\gamma,o,j,c,t,p_1,p_2,p_3,\xi.
\]
Here $z$ remains the transport coordinate and $t$ is the joint-guard
root coordinate. Set $\lambda=2^{1113}$ and $\mu=2^{666104}$; the other
four fixed coefficients are the complete source34 tuple
$\beta,K_0,M_C,M_F$, with $\beta=2^{1064921}-1$ and $M_F$ already shifted.
Use only the following polynomial abbreviations:
\begin{align*}
 J&=3J_3+1,&q&=\beta J+1,&W&=4(3w_3+1),\\
 X&=qW,&Y&=(2v+1)q^3,&E&=XY,\\
 a&=(X+1)Y,&H&=4a+3,&\Delta&=a^2+H,\\
 V&=fo-c,&k&=1-V+jc+hE,&\eta&=c-Yk,\\
 u&=EYk,&A&=u(u+k)+1,&I&=(2r+1)c^2,\\
 Q&=I^2,&\delta&=Q-1,&F&=(K_0+W)C-z(q-1).
\end{align*}
Further put
\begin{align*}
 R&=\bigl(q(q-F)-C\bigr)(q^2-1)+(M_C+qM_F)J,\\
 d&=(\lambda F)(\mu C)-(q-\lambda F-\mu C)^2,\\
 T&=\eta(k-\eta),\\
 P_{\rm pos}&=1+p_1^2+p_2^2+p_3(p_3+1),\\
 G&=dT-(t^2-d-T+1)^2-P_{\rm pos},\\
 B_0&=QV^2-1,\\
 L&=\xi^2+2I(2u+k)\xi+Q(k^2-4),\\
 D_n&=\delta L-4B_0,\\
 N&=D_n^2-64\delta QAB_0.
\end{align*}
The promised single equation is
\begin{equation}\label{eq:full17}
 \boxed{P_{17}=M^2+B_p^2+S_p^2+G^2+N^2=0,}
\end{equation}
where
\[
 M=(X+ac+\gamma H)^2-\Delta c^2-1,\quad
 B_p=cj-V-R,\quad S_p=Q-\Delta(f^2-1).
\]
There are no hidden quantified coordinates or divisions in this display.
Its five residual degrees are $14,5,14,44,56$. The coefficient of
\[
 r^{12}c^{24}h^4f^4o^4J_3^{44}w_3^8v^{12}
\]
in $P_{17}$ is exactly $2^{52}3^{52}\beta^{44}>0$, certifying degree
$112$ after specialization to the actual coefficient tuple. The supplied
literal arithmetic DAG uses $124$ binary operations under the stated
counting convention; this is a verified construction, with no claim of
operation minimality.

The scale $I$ in the merger is essential. Taking $I=2$, $Q=4$, $V=2$,
and $A=5$ makes the unscaled auxiliary radicand
$(QV^2-1)/(Q-1)=5$. The unscaled norm then accepts the spurious rational
sum $\sqrt5-\sqrt5=0$, although neither square root is integral. The
scaled radicands are $20$ and $5$, and the same sum-zero input is rejected.
